% Generated from ../chapters/16-an-unplanned-twenty-year-experiment.md
\chapter{An Unplanned Twenty-Year Experiment}

My practice did not begin as a longevity study. Tai Chi, Qigong, and meditation were pursued for their immediate effects and because the traditions themselves emphasized relaxation, alignment, attention, and self-cultivation.

During the early years, I noticed changes that were difficult to describe using the ordinary categories of exercise. The body sometimes appeared to adjust itself when conscious effort decreased. Areas of tension released in sequences. Breathing reorganized. The relation among head, jaw, neck, spine, and pelvis changed. Audible events occurred during deeply relaxed practice.

Over years, the pattern seemed less like acquiring a new shape and more like removing compensations.

\section{What is actually observed}

The strongest defensible observations are phenomenological and functional:

\begin{itemize}
\tightlist
\item
  repeated, reproducible sensations of release and reorientation;
\item
  changes in habitual posture and movement;
\item
  altered muscular effort;
\item
  changes in breathing and balance;
\item
  audible cervical or temporomandibular events;
\item
  perceived reduction of asymmetry;
\item
  persistence of change over long periods.
\end{itemize}

These observations do not identify anatomical mechanisms. Sound cannot distinguish cavitation, tendon movement, disc displacement, or bone remodelling. Appearance cannot distinguish posture, muscle, fat, fluid, dental change, or skeletal geometry.

The correct response is neither dismissal nor premature certainty. It is measurement.

\section{Why personal observation matters}

An N-of-1 observation is weak evidence for general efficacy but valuable for hypothesis generation, especially when the timescale is longer than typical studies.

Scientific history often begins with an anomaly that does not fit the available model. The anomaly earns attention when it is persistent, reproducible, and capable of generating discriminating predictions.

The value of my account therefore depends on what follows:

\begin{itemize}
\tightlist
\item
  standardized photographs and 3D scans;
\item
  range-of-motion and force measures;
\item
  electromyography;
\item
  dental and temporomandibular assessment;
\item
  clinically justified imaging;
\item
  gait and balance analysis;
\item
  longitudinal biomarkers;
\item
  independent replication.
\end{itemize}

\section{The danger of narrative}

Twenty years of practice also create opportunities for confirmation bias. Memory reconstructs. Aging, dental work, body composition, camera angle, and unrelated health changes can be misattributed.

A serious framework must make it easier to discover that the interpretation is wrong.

The personal story is not presented as proof of endogenous rejuvenation. It is the source of the questions that organize this book.
